\begin{scanpage}{0214}{196}
\[\left.\begin{aligned}
\mathcal E^{(n)}(z_1)&=-\mathcal E_2^{(n-1)}\psi_-(L)e^{ikL},\\
\mathcal E^{(n)}(z_2)&=-\mathcal E_1^{(n-1)}\psi_+(L)e^{ikL}
\end{aligned}\right\}\tag{31.15}\]
and
\[\mathcal E_1^{(n)}=\left.\mathcal E^{(n)}(z_1)\right|_{\vartheta=0},\quad\mathcal E_2^{(n)}=\left.\mathcal E^{(n)}(z_2)\right|_{\vartheta=\pi}.\tag{31.16}\]
Thus, the field arising with multiple diffractions (starting with the second) may be represented in the following form:
\[\begin{aligned}
&\sum_{n=2}^\infty\bigl[E_\vartheta^{(n)}(z_1)+E_\vartheta^{(n)}(z_2)\bigr]\\
&\quad=\frac{1}{2\sin\vartheta\ln\dfrac{2i}{\gamma ka\sin\vartheta}}\frac{e^{ikR}}R
\left[\sum_{n=2}^\infty\mathcal E^{(n)}(z_1)e^{-ikz_1\cos\vartheta}
+\sum_{n=2}^\infty\mathcal E^{(n)}(z_2)e^{-ikz_2\cos\vartheta}\right],
\end{aligned}\tag{31.17}\]
where
\[\sum_{n=2}^\infty\mathcal E^{(n)}(z_1)=\mathcal E^{(2)}(z_1)+\bigl[\mathcal E_1^{(2)}\psi(L)e^{ikL}-\mathcal E_2^{(2)}\bigr]\frac{\psi_-(L)}D e^{ikL},\tag{31.18}\]
\[\sum_{n=2}^\infty\mathcal E^{(n)}(z_2)=\mathcal E^{(2)}(z_2)+\bigl[\mathcal E_2^{(2)}\psi(L)e^{ikL}-\mathcal E_1^{(2)}\bigr]\frac{\psi_+(L)}D e^{ikL},\tag{31.19}\]
and the functions $\mathcal E^{(2)}(z_{1,2})$ and $\mathcal E_{1,2}^{(2)}$ are determined by Equations (30.31), (30.35), (31.02) and (31.04). We will not write out here the rather unwieldy final expression for this field, but we will proceed with a calculation of the total field scattered by a vibrator.
\sectiontitle{32}{Total Fringing Field}
Before beginning the derivation of the expression for the scattering characteristic, let us make the following observation. The
\end{scanpage}
